\paragraph{The robust part of Stein's theorem.} In general relativity, under
the premises of Remark~\ref{rem:lc}, nothing outside an event's causal past
can have become for it unless everything has. This holds in every globally
hyperbolic spacetime, however asymmetric, and does not rely on any foliation.
Those who want an extended, spacelike ``now'' as part of objective becoming
face Stein's dilemma in general relativity too, unless they change the
category or reject locality on the metric-only category.

\paragraph{The strongest rival: clock-equipped spacetimes.} Let $\GH_T$ have as
objects pairs $(M,T)$ with $M\in\GH$ and $T$ a smooth Cauchy temporal function, and
as morphisms $\psi:(M,T)\to(M',T')$ the morphisms of $\GH$ with $T'\circ\psi=T$. The
rule $x\,B_{(M,T)}\,b\iff T(x)\le T(b)$ is transitive, contains $\Ichr$, satisfies
(LC) on this category, and has the level sets of $T$ as shared presents. The
Milne inclusion of Proposition~\ref{prop:nonlocal} is not a morphism here: its
clock $\tau_W=(t^2-|\mathbf x|^2)^{1/2}$ has no smooth extension across the light
cone, where its gradient diverges. So this rival is local with respect to
its own morphisms. It rejects one premise, that becoming is a function of the
metric alone over all of $\GH$, and Theorem~\ref{thm:main} then says nothing about
it (in effect Monton's move). Its cost depends on the clock. A clock definable
from the metric, as cosmological time is in big-bang cosmologies, adds no
structure; the price is that the rule gives no verdict where that clock does
not exist, and locality no longer selects it, since $\Jm$ also satisfies the
axioms on $\GH_T$. A preferred foliation or reference field not definable
from the metric is added structure. Rejecting locality on $\GH$ itself is a
distinct option (Proposition~\ref{prop:nonlocal}).

\paragraph{A trilemma.} Suppose a shared present is understood as a set of
distinct events each of which has become for each other, the sense in which
Stein's dilemma is posed. Then Corollary~\ref{cor:nopresent} can be read as a
trilemma for becoming theories, which are transitive by definition, defined from the metric alone: locality, a
shared (extended) present, and non-trivial becoming cannot all be had. Defenders of a shared now reject locality on $\GH$, reject transitivity or change the category; local becoming theorists get
point-like presents, which recalls Stein's remark that the converse of
Putnam's premise leads to ``a peculiarly extreme (but pluralistic!) form of
solipsism'' \citep[p.~18]{Stein1968}; and giving up non-triviality is the
block. Other notions of a present, for instance a set of events none of which
has become for another, are not addressed by the corollary. The corollary
says nothing about which horn to take.

\paragraph{The fragile part.} Uniqueness does not survive. Minkowskian
symmetry forced a choice between $\Ichr$ and $\Jm$; general relativity leaves
open how lightlike-separated events stand, and a rule may let curvature along
the light ray decide. Clifton and Hogarth's debate over causal versus
chronological becoming \citep[pp.~364--365]{CH1995} thus acquires a third
family of options; the choice between the endpoints is equivalent to a
regularity requirement (Proposition~\ref{prop:reg}).

\paragraph{Locality does the work symmetry did.} In a single curved spacetime,
definability-style invariance constrains almost nothing
(Proposition~\ref{prop:single}); natural global rules such as cosmological
time violate only locality (Proposition~\ref{prop:nonlocal}). So the
curved-spacetime case for ``only the past has become'' rests on a locality
principle of the kind locally covariant field theory takes for granted and
akin to Viglione's Causal Connection Condition \citep{Viglione2026} (which is
stricter: (LC) allows dependence on spacelike events inside the convex region). Rejecting
locality is a coherent option, but it is the option of global, cosmological
becoming, with its own costs \citep{Dieks2006,Viglione2026}.

\paragraph{Limits.} The theorem concerns relations of becoming, not
ontology; it does not decide between presentism, the growing block and
eternalism. An eternalist can accept $\Jm$ as what has happened by an event.
Reading $B_M$ as existence from $b$ is a further bridge, which nothing here
supplies. The theorem assumes that becoming is defined for every globally hyperbolic
geometry (Remark~\ref{rem:offshell}); Section~\ref{sec:energy} shows that
dropping this premise is not routine. Spacetimes that are not globally hyperbolic are not treated beyond
Proposition~\ref{prop:godel}. The time orientation is an input, through (CB);
the theorem shows that, given an orientation, only the past lobe can have
become, but it does not explain the orientation. Nor does it address quantum
nonlocality. Clifton and Hogarth's maxim that whatever can causally influence
an event has become for it \citep[p.~364]{CH1995} needs care. Read as a claim
about signalling, it puts $\Jm$ into $B_M$. Bell-type dependence on spacelike
events is not signalling, so it is no counterexample to that reading. What it
challenges is the stronger thought that any dynamical dependence on an event
requires that event to have become. If the fundamental dynamics lets spacelike
events contribute to bringing an event about, the stronger reading puts them
into $B_M$, and Theorem~\ref{thm:main} makes becoming universal. A non-trivial
local becoming theory that holds the stronger reading therefore presupposes that
such contribution is restricted to the causal past. That is a bet about the
dynamics, not a theorem of relativity.
